%% Chapter 5 — Implementation
\chapter{Implementation}
\label{ch:implementation}

This chapter documents the technical realisation of the NPR avatar system. It covers the development environment, the integration path for each avatar platform, the implementation of the shared rendering structure, and the HLSL-level details of each stylisation technique developed to date. Runtime tooling developed to support iterative development is described at the end of the chapter.

\section{Development Environment}
\label{sec:impl_environment}

The system was developed in Unity~6 (version 6000.0.60f1) using the Universal Render Pipeline~(URP). The target hardware is the Meta Quest family of standalone VR headsets. The Meta Avatars SDK version~40.0.1 was used for the SDK-managed avatar platform. All shader code is written in HLSL and targets the URP forward rendering path.

\section{Avatar Platform Integration}
\label{sec:impl_platforms}

\subsection{Mixamo}
\label{sec:impl_mixamo}

Mixamo characters are exported as FBX assets and imported into Unity as standard skinned mesh renderers. No proprietary shader system is involved, so all URP materials can be replaced directly with custom NPR shaders without restriction. The Mixamo character Jody served as the primary prototyping target during technique development, allowing the NPR shaders to be iterated freely before the more constrained platform integrations were addressed.

\subsection{Meta Avatars SDK}
\label{sec:impl_meta}

The Meta Avatars SDK manages avatar materials through its Extensible Shader Framework~(ESF), which generates platform-specific HLSL from parameterised templates. The generated source is not intended to be edited directly. Custom logic is instead injected through named hook functions that the ESF exposes at defined points in the shader pipeline.

The hook used in this thesis is \texttt{AppSpecificPostManipulation}, which executes per fragment after all SDK PBR lighting — subsurface scattering, eye glints, hair shading — has been computed. It receives the accumulated lit colour and returns a modified colour, giving full control over the final fragment output without disrupting any SDK rendering logic. The hook is defined across three files that the ESF includes automatically during shader compilation:

\begin{itemize}
  \item \texttt{app\_functions.hlsl} — the body of \texttt{AppSpecificPostManipulation}; dispatches to the active NPR technique based on the \texttt{\_NPREffect} keyword
  \item \texttt{app\_declarations.hlsl} — uniform declarations for all NPR parameters, shared across techniques
  \item \texttt{app\_variants.hlsl} — shader feature keyword declarations that activate each technique branch
\end{itemize}

A further complication is the SDK's LOD management system, which replaces avatar materials automatically as camera distance changes. This reverts any custom shader assignment made at load time. The \texttt{AvatarShaderSwapper.cs} script addresses this by running a periodic coroutine that re-applies the custom NPR materials after each LOD transition is detected. The script also handles the initial material replacement, waiting for the asynchronous avatar load to complete before copying base colour and normal textures from the SDK materials into the custom shader properties.

\subsection{Avaturn}
\label{sec:impl_avaturn}

Avaturn exports avatars in the \texttt{.glb} format, the binary container of the glTF~2.0 standard. Unity does not natively support runtime glTF loading, so the GLTFast package is used to load the asset at runtime. On load completion, GLTFast produces a standard Unity GameObject hierarchy with skinned mesh renderers and PBR materials assigned from the glTF material definitions, including albedo, normal, and metallic-roughness textures. Because Avaturn avatars carry no proprietary shader constraints once imported, the custom NPR materials are assigned to the resulting skinned meshes using the same approach as Mixamo: direct material replacement after load. Blendshape animation targets exported by Avaturn are preserved through the GLTFast import and driven at runtime via \texttt{SkinnedMeshRenderer.SetBlendShapeWeight} for the video recording scenarios used in the user study.

\section{Shared Rendering Structure}
\label{sec:impl_shared}

The NPR shader used across Mixamo and Avaturn is \texttt{Avatar-Meta-UGB.shader}, a multi-pass URP shader containing two passes: \texttt{NPROutline} and \texttt{ForwardLit}. For the Meta SDK platform, the forward-lit NPR logic runs inside \texttt{AppSpecificPostManipulation}; the outline pass is handled by a separate material applied to a duplicate mesh layer.

The core forward-lit dispatch lives in \texttt{Style2MetaAvatarCore.hlsl}. It reads the active \texttt{\_NPREffect} shader keyword and branches to the corresponding technique include file. This structure means all eight techniques share the same entry point and parameter infrastructure; switching techniques at runtime requires only toggling the active keyword via \texttt{Material.EnableKeyword} and \texttt{Material.DisableKeyword}, which the runtime UI performs without recompilation.

\subsection{Inverted Hull Outline Pass}
\label{sec:impl_outline}

The outline pass implements the inverted hull silhouette. In the Mixamo and Avaturn shaders the pass sets \texttt{Cull Front}, \texttt{ZWrite On}, \texttt{ZTest Less}; in the Meta SDK shader it sets \texttt{Cull Front}, \texttt{ZWrite Off}, \texttt{ZTest LEqual} to avoid conflicting with the SDK's depth compositing. The vertex shader displaces each vertex outward along its world-space normal by \texttt{\_OuterOutlineWidth} (Mixamo/Avaturn, world-space metres) or \texttt{\_OutlineWidth * 0.001} (Meta, converted from UI units). The fragment outputs a flat \texttt{\_OuterOutlineColor} or \texttt{\_OutlineColor}. An \texttt{\_OutlineDepthBias} property shifts the clip-space Z coordinate toward the camera to prevent z-fighting where the expanded shell and the original surface meet at shallow angles.

\subsection{Alpha Cutout Handling}
\label{sec:impl_alpha}

Alpha cutout is controlled by the \texttt{\_EnableAlphaTest} flag rather than being always-on. When enabled, both the outline pass and the forward-lit pass sample the base colour texture alpha channel and call \texttt{clip(alpha - threshold)} before any displacement or shading is applied. In the outline pass this ensures the expanded hull is discarded wherever the hair card texture is transparent, preventing a rectangular border from appearing around each card. The cutout threshold is exposed in the runtime UI and adjusted per avatar submesh — a lower threshold is typically used for eyelashes (dense geometry) than for hair cards (sparse, large transparent regions).

\section{NPR Technique Implementation}
\label{sec:impl_techniques}

\subsection{V1 — Toon Shading}
\label{sec:impl_v1}

Three separate shader implementations exist for V1, one per avatar platform. The outline pass is structurally identical across all three (inverted hull, normal-offset displacement, alpha cutout); the toon shading pass differs because the Mixamo and Avaturn variants operate on raw NdotL, while the Meta variant operates on the composited PBR output.

\subsubsection{Mixamo}
\label{sec:impl_v1_mixamo}

\textbf{Shader:} \texttt{Assets/AvatarShaderExperimental/Shaders/V1\_ToonShading\_GeometryOutline.shader} \\
\textbf{Pipeline target:} URP, \texttt{\#pragma target 3.0}

The outline pass sets \texttt{Cull Front}, \texttt{ZWrite On}, \texttt{ZTest Less} and displaces each vertex along its world-space normal by \texttt{\_OuterOutlineWidth} (world units). The fragment returns a flat \texttt{\_OuterOutlineColor}. When \texttt{\_EnableAlphaTest} is active, a \texttt{clip} call discards fragments below the alpha threshold before any displacement is applied, ensuring the hull respects hair card silhouettes. An \texttt{\_OutlineDepthBias} value shifts clip-space Z toward the camera to prevent z-fighting at shallow angles.

The toon pass quantises the Lambertian dot product $\mathbf{N} \cdot \mathbf{L}$ per step using a per-band smoothstep. The expression after unification with the Avaturn version is:

\begin{lstlisting}[language=C, style=mystyle, caption={V1 toon quantisation — Mixamo and Avaturn (unified).}]
float steps  = max(1.0, _ToonSteps);
float scaled = NdotL * steps;
float band   = floor(scaled);
float frac   = scaled - band;
float blend  = smoothstep(1.0 - _ToonSmoothness, 1.0, frac);
float toon   = saturate((band + blend) / steps);
toon = lerp(1.0 - _ShadowStrength, 1.0, toon);
\end{lstlisting}

Each band boundary receives an independent narrow blend of width \texttt{\_ToonSmoothness}, so normal jitter near any seam produces a localised crossfade rather than snapping the whole surface patch between bands. Vertex positions and normals are read directly from the standard mesh attributes; no compute-skinning bridge is needed for standard \texttt{Skinned\-Mesh\-Renderer} avatars. An additive ambient \texttt{\_AmbientColor} and a rim term \texttt{pow(1 - NdotV, RimPower)} are added to the quantised diffuse.

\subsubsection{Avaturn}
\label{sec:impl_v1_avaturn}

\textbf{Shader:} \texttt{Assets/Shaders/NPR/V1\_ToonShading\_GeometryOutline.shader} \\
\textbf{Pipeline target:} URP, \texttt{\#pragma target 3.5} \\
\textbf{Materials:} five materials covering body, head, hair, eyelash, and look submeshes

The toon pass uses the same per-step quantisation as Mixamo (Listing~5.1). The outline pass is structurally identical to Mixamo with one critical addition: before computing world positions, the vertex shader calls \texttt{OVR\_FETCH\_POS\_NORM(v.vertex.xyz, v.normal, v.vertexID)}. This bridge reads vertex positions and normals from the external compute-skinned vertex buffers written by the Meta SDK runtime. Without it, the GPU-skinned vertices are not visible to the vertex shader, and the outline hull renders at the bind-pose position. The same bridge is required in the toon pass for correct deformation tracking. For standard Avaturn avatars loaded via GLTFast without the Meta SDK runtime, the macro resolves to a no-op and the attributes fall back to the standard mesh inputs.

\subsubsection{Meta Avatars SDK}
\label{sec:impl_v1_meta}

\textbf{Shader:} \texttt{Assets/Shaders/CustomShaders/Avatar-Meta-UGB.shader} \\
\textbf{Keyword:} \texttt{EFFECT\_TOON} (enabled via \texttt{ENABLE\_NPR\_EDGES}) \\
\textbf{Technique file:} \texttt{Assets/Shaders/CustomShaders/NPREffect\_Toon.cginc}

The Meta implementation runs inside \texttt{AppSpecificPostManipulation}, which receives the fully composited PBR colour produced by the SDK's \texttt{STYLE\_2\_STANDARD} renderer. The quantisation therefore operates on RGB directly, preserving hue while stepping luminance:

\begin{lstlisting}[language=C, style=mystyle, caption={V1 posterisation — Meta SDK (\texttt{NPREffect\_Toon.cginc}).}]
float  bands      = max(2.0, _ToonColorBands);
float3 posterized = floor(color.rgb * bands + 0.5) / bands;
color.rgb = lerp(color.rgb, posterized, _ToonPosterizeStrength);

float lum  = dot(color.rgb, float3(0.2126, 0.7152, 0.0722));
color.rgb  = lerp(float3(lum, lum, lum), color.rgb, _ToonSaturation);
\end{lstlisting}

The \texttt{floor(rgb * bands + 0.5) / bands} expression rounds each channel to the nearest multiple of $1/\text{bands}$ (round-to-nearest rather than floor-to-band), which distributes quantisation error symmetrically and avoids the darkening bias of a plain floor. Saturation adjustment runs after posterisation to prevent clamping artefacts. Because the input already contains Meta's subsurface scattering, anisotropic hair, and eye glints, the posterised output stylises a physically correct image rather than reconstructing lighting from scratch.

The outline pass in this shader is a dedicated \texttt{NPROutline} pass with \texttt{Cull Front}, \texttt{ZWrite Off}, \texttt{ZTest LEqual}. The vertex shader displaces each SDK-skinned vertex outward by:

\begin{lstlisting}[language=C, style=mystyle, caption={Outline vertex displacement — Meta SDK outline pass.}]
float3 worldPos = o.v_WorldPos + normalize(o.v_Normal) * _OutlineWidth * 0.001;
\end{lstlisting}

The \texttt{0.001} scale converts \texttt{\_OutlineWidth} from millimetre-scale UI units to metre-scale world space. The outline is toggled via \texttt{\_OutlineEnabled}; the fragment discards if the value is below 0.5. Unlike the Mixamo and Avaturn variants, \texttt{ZWrite Off} is used because the Meta pipeline already has correct depth from the main pass and a second depth write would cause artefacts with the SDK's transparency compositing.

\begin{table}[H]
  \centering
  \caption{V1 implementation comparison across the three avatar platforms.}
  \label{tab:v1_comparison}
  \small
  \begin{tabularx}{\textwidth}{lXXX}
    \toprule
    & \textbf{Mixamo} & \textbf{Avaturn} & \textbf{Meta SDK} \\
    \midrule
    Quantisation method & Per-step smoothstep, NdotL & Per-step smoothstep, NdotL & Round-to-nearest, composited RGB \\
    Skinning bridge & None & \texttt{OVR\_FETCH\_POS\_NORM} & SDK native \\
    Rim light & Additive, computed & Additive, computed & Pre-composited by Meta PBR \\
    Outline depth write & \texttt{ZWrite On} & \texttt{ZWrite On} & \texttt{ZWrite Off} \\
    Outline toggleable & No & No & Yes (\texttt{\_OutlineEnabled}) \\
    Shadow source & URP \texttt{GetMainLight} & URP cascades & Meta PBR attenuation \\
    \bottomrule
  \end{tabularx}
\end{table}

\begin{figure}[H]
  \centering
  \missingfigure[figwidth=\textwidth]{V1 (Toon Shading): Mixamo (left) / Meta SDK (centre) / Avaturn (right)}
  \caption{V1 (Toon Shading) across the three avatar platforms.}
  \label{fig:v1_impl}
\end{figure}

\subsection{V2 — X-Toon Extended Toon Shading}
\label{sec:impl_v2}

The X-Toon implementation replaces the procedural quantization of V1 with a lookup into a two-dimensional ramp texture, \texttt{\_ToonRampTex}. The horizontal axis of the ramp is indexed by the Lambertian dot product $\mathbf{N} \cdot \mathbf{L}$, remapped from $[-1, 1]$ to $[0, 1]$. The vertical axis is indexed by a depth factor derived from the linear eye-space depth of the fragment, normalised against a configurable \texttt{\_DepthRange} parameter. Sampling the ramp with these two coordinates produces a colour that encodes both the diffuse shading band and its depth-dependent modification. The Fresnel rim and ambient terms from V1 are retained. The ramp texture itself is authored externally and assigned as a material property, giving the designer full control over the shading character without shader recompilation.

\begin{figure}[H]
  \centering
  \missingfigure[figwidth=\textwidth]{V2 (X-Toon): Mixamo (left) / Meta SDK (centre) / Avaturn (right)}
  \caption{V2 (X-Toon Extended Toon Shading) across the three avatar platforms.}
  \label{fig:v2_impl}
\end{figure}

\subsection{V3 — Screen-Space Normal Edge Detection}
\label{sec:impl_v3}

The normal edge implementation uses the \texttt{ddx} and \texttt{ddy} hardware intrinsics applied to the interpolated world-space surface normal $\mathbf{N}$, which is available as a fragment-stage interpolant at no additional texture sample cost. The gradient magnitude is computed as $M_N = \sqrt{\|\texttt{ddx}(\mathbf{N})\|^2 + \|\texttt{ddy}(\mathbf{N})\|^2}$, producing a scalar per fragment that is high at surface creases, mesh boundaries, and regions of high curvature. This value is compared against \texttt{\_NormalThreshold}; fragments exceeding the threshold are output as \texttt{\_EdgeColor}, and all others pass the base albedo through unchanged. An optional \texttt{\_NormalEdgePower} exponent controls the response curve, allowing soft crease edges to be suppressed relative to sharp geometric boundaries.

\begin{figure}[H]
  \centering
  \missingfigure[figwidth=\textwidth]{V3 (Normal Edge Detection): Mixamo (left) / Meta SDK (centre) / Avaturn (right)}
  \caption{V3 (Screen-Space Normal Edge Detection) across the three avatar platforms.}
  \label{fig:v3_impl}
\end{figure}

\subsection{V4 — Sobel Edge Detection}
\label{sec:impl_v4}

The Sobel implementation samples the diffuse albedo texture at eight neighbours in a $3\times3$ grid centred on the current fragment, using a fixed texel offset \texttt{\_TexelSize}. The horizontal and vertical gradient components $G_x$ and $G_y$ are accumulated using the Sobel kernel weights, and the edge magnitude is computed as $M = \sqrt{G_x^2 + G_y^2}$.

Two artefact suppression mechanisms are applied before the threshold test. The luminance range check computes the spread of luminance values across the nine samples; if the spread exceeds \texttt{\_SeamRangeLimit}, the pixel is rejected as a UV seam artefact. The skin classification converts the centre sample to HSV and compares the saturation channel against \texttt{\_SkinSaturationCutoff}; fragments below the cutoff are classified as skin and tested against a stricter threshold \texttt{\_SkinEdgeThreshold} rather than the standard \texttt{\_EdgeThreshold}. Fragments exceeding the active threshold are coloured \texttt{\_EdgeColor}; all others pass through the base albedo.

\begin{figure}[H]
  \centering
  \missingfigure[figwidth=\textwidth]{V4 (Sobel): Mixamo (left) / Meta SDK (centre) / Avaturn (right)}
  \caption{V4 (Sobel Edge Detection) across the three avatar platforms.}
  \label{fig:v4_impl}
\end{figure}

\subsection{V5 — Gaussian Prefiltered Sobel}
\label{sec:impl_v5}

The Gaussian prefiltered Sobel implementation adds a smoothing stage before gradient computation. A $5\times5$ sample grid is evaluated at the current fragment, with each sample weighted by a precomputed Gaussian coefficient derived from a standard deviation \texttt{\_GaussianSigma}. The weighted sum produces a smoothed colour estimate at each of the nine Sobel positions. The Sobel gradient is then applied to the smoothed values rather than the raw texture samples. Because the Gaussian and Sobel stages both require neighbourhood samples, the total sample count for this technique is higher than V4; the exact count depends on the kernel radius configured via \texttt{\_GaussRadius}. The artefact suppression mechanisms from V4 are retained.

\begin{figure}[H]
  \centering
  \missingfigure[figwidth=\textwidth]{V5 (Gaussian Sobel): Mixamo (left) / Meta SDK (centre) / Avaturn (right)}
  \caption{V5 (Gaussian Prefiltered Sobel) across the three avatar platforms.}
  \label{fig:v5_impl}
\end{figure}

\subsection{V6 — Hierarchical Edge Detection}
\label{sec:impl_v6}

The hierarchical implementation computes three independent edge layers in a single fragment shader invocation and combines them via weighted addition.

The depth layer uses \texttt{ddx} and \texttt{ddy} applied to the fragment's linear eye depth to estimate the depth gradient magnitude. The gradient is normalised by \texttt{\_DepthSensitivity} and thresholded to produce a binary edge signal. The normal layer applies the same derivative approach to the interpolated world-space surface normal vector, computing $M_N = \sqrt{\|\texttt{ddx}(\mathbf{N})\|^2 + \|\texttt{ddy}(\mathbf{N})\|^2}$ and thresholding against \texttt{\_NormalThreshold}. The colour layer applies the Roberts Cross operator to the diffuse texture: it samples four neighbours in a $2\times2$ diagonal pattern and computes the cross-difference magnitudes $|R_1|$ and $|R_2|$ as defined in~\Cref{sec:background_gradient}, thresholded against \texttt{\_ColorThreshold}.

The three binary edge signals are weighted by \texttt{\_DepthWeight}, \texttt{\_NormalWeight}, and \texttt{\_ColorWeight} respectively and summed. The combined weight is clamped to $[0,1]$ and used to blend between the base albedo and \texttt{\_EdgeColor}.

\begin{figure}[H]
  \centering
  \missingfigure[figwidth=\textwidth]{V6 (Hierarchical): Mixamo (left) / Meta SDK (centre) / Avaturn (right)}
  \caption{V6 (Hierarchical Edge Detection) across the three avatar platforms.}
  \label{fig:v6_impl}
\end{figure}

\subsection{V7 — Kuwahara Painterly Filter}
\label{sec:impl_v7}

The Kuwahara implementation evaluates four overlapping quadrant subregions around the current fragment. Each quadrant is a square of radius \texttt{\_KuwaharaRadius}, offset from the centre so that adjacent quadrants share a one-texel border. For each quadrant, the implementation accumulates the sum and sum-of-squares of the sampled colour values, then computes the per-channel mean and variance from these accumulators. The quadrant with the lowest luminance variance is selected, and its mean colour is written as the fragment output.

Because the filter operates on the diffuse texture using fixed texture-coordinate offsets, no intermediate render target is required, and the entire computation runs inside the single-pass fragment hook. The isotropic formulation is used; the anisotropic extension of Kyprianidis et al.~\cite{kyprianidis2009} requires a structure tensor computed in a prior pass and is therefore incompatible with the single-pass constraint of the Meta SDK hook.

\begin{figure}[H]
  \centering
  \missingfigure[figwidth=\textwidth]{V7 (Kuwahara): Mixamo (left) / Meta SDK (centre) / Avaturn (right)}
  \caption{V7 (Kuwahara Painterly Filter) across the three avatar platforms.}
  \label{fig:v7_impl}
\end{figure}

\subsection{V8 — Kuwahara with Sobel Edges}
\label{sec:impl_v8}

The V8 implementation chains the Kuwahara and Sobel stages in sequence within the same fragment shader invocation. The Kuwahara filter is applied first, producing a smoothed mean colour from the lowest-variance quadrant. The Sobel gradient is then computed using the Kuwahara output as the input signal rather than the raw diffuse texture. Because the Kuwahara output is piecewise uniform within regions, the Sobel operator responds primarily to inter-region boundaries, suppressing the fine texture gradients that produce noise in V4. The skin saturation and seam rejection checks from V4 are retained. The detected edges are composited over the Kuwahara output using \texttt{\_EdgeColor} and \texttt{\_EdgeThreshold}.

\begin{figure}[H]
  \centering
  \missingfigure[figwidth=\textwidth]{V8 (Kuwahara + Sobel): Mixamo (left) / Meta SDK (centre) / Avaturn (right)}
  \caption{V8 (Kuwahara with Sobel Edges) across the three avatar platforms.}
  \label{fig:v8_impl}
\end{figure}

\subsection{V9 — Composite: Kuwahara, Gaussian Sobel, and Hierarchical}
\label{sec:impl_v9}

The composite technique chains three stages. First, the Kuwahara filter produces the painterly base. Second, the Gaussian prefiltered Sobel is applied to the Kuwahara output to extract smoothed colour boundary lines. Third, the hierarchical depth, normal, and Roberts Cross edge layers are evaluated and composited on top of the Gaussian Sobel result with independent weighting. All three stages share the same input texture and are computed sequentially within the single fragment shader invocation. The total sample count is the sum of the Kuwahara quadrant samples, the Gaussian prefilter samples, and the four Roberts Cross samples, making this the most sample-intensive technique developed to date.

\begin{figure}[H]
  \centering
  \missingfigure[figwidth=\textwidth]{V9 (Composite): Mixamo (left) / Meta SDK (centre) / Avaturn (right)}
  \caption{V9 (Composite: Kuwahara, Gaussian Sobel, Hierarchical) across the three avatar platforms.}
  \label{fig:v9_impl}
\end{figure}

\section{Runtime Tooling}
\label{sec:impl_tooling}

\subsection{In-VR Parameter Tuning Interface}
\label{sec:impl_ui}

\texttt{NPREdgeDetectionUI.cs} implements a fully procedural WorldSpace UI panel that is visible and interactive from inside the headset. The panel is constructed at runtime from Unity UI primitives: each tunable parameter is represented by a labelled row containing either a slider or a toggle button, laid out in technique-specific groups. Controller raycasting using the OVR interaction layer detects when the ray from either controller intersects a slider's \texttt{BoxCollider}; a trigger-pull gesture initiates a drag interaction that maps horizontal ray movement to the slider's value range. Technique switching is handled through a dropdown row: selecting a technique calls \texttt{EnableKeyword} and \texttt{DisableKeyword} on the active materials to swap the active shader branch without recompilation.

The panel exposes parameters across all technique groups, including per-technique thresholds, weights, kernel radii, and colour properties. This tool was central to the visual assessment process used to select the technique for the user study.

\subsection{LOD Material Management}
\label{sec:impl_swapper}

\texttt{AvatarShaderSwapper.cs} solves the material reversion problem described in~\Cref{sec:impl_meta}. On startup it subscribes to the avatar load completion callback and performs an initial material replacement pass across all LOD levels. A coroutine running at a configurable polling interval inspects the active materials on each renderer at runtime; if an SDK-assigned material is detected — indicating that a LOD transition has reverted the custom assignment — the coroutine re-applies the NPR material and copies the base colour and normal texture references from the newly active SDK material into the replacement shader's properties.

\subsection{Pose Freeze for Evaluation}
\label{sec:impl_freeze}

\texttt{AvatarFreezeController.cs} enables the avatar's body-tracking pose to be frozen at a chosen instant, preserving a specific expression and body posture for evaluation screenshots. The freeze is implemented by accessing the \texttt{\_inputTrackingProvider} field of \texttt{OvrAvatarInputManager} via reflection and substituting a snapshot provider that returns a fixed pose state. Pressing A, X, or the Space key toggles the freeze on and off.
